\chapter{Implications}
\label{ch:essay-implications}

The theorems proven in \emph{Admissible Reconstruction} establish consequences that follow necessarily from the admissibility framework and generate hypotheses that require independent empirical validation. The distinction between these two categories must be maintained throughout.

\section{Established Consequences}

Theorem~4.2, the structural direction of Non-Identifiability proven in Chapter~4, establishes that observation maps constrained to preserve only reachable-set structure are generically non-injective whenever the space of histories exceeds what that structure can encode. This applies directly to historical linguistics. Sound change operating within a language lineage and lexical borrowing from a neighboring language can produce phonologically identical outcomes in the descendant language despite originating from entirely different historical processes. If two evolutionary histories $H_1$ and $H_2$, one involving regular sound change and the other involving contact-induced borrowing, both lead to the same phonological form in the observable contemporary language, and if that phonological form determines the same set of admissible continuations regardless of which history produced it, then the two histories share the same reachable-set structure: $\reach{H_1} = \reach{H_2}$. Under these conditions, Theorem~4.2 guarantees that any observation map $\obsmap$ factoring through reachable-set structure produces a collision: $\obsmap(H_1) = \obsmap(H_2)$. No reconstruction procedure operating solely on the observable contemporary form can determine with certainty which history actually occurred.

The diagnostic corollary, Corollary~4.1, provides a criterion for assessing such cases. An observer examining a contemporary form and suspecting that it could have arisen either through inheritance with regular sound change or through borrowing from a contact language can, in principle, check whether the two candidate histories would produce the same reachable set. If the reachable sets differ, then additional evidence—perhaps from comparative data, geographical distribution, or semantic fields—could distinguish the histories, and the reconstruction uncertainty is correctable. If the reachable sets are equal, then the uncertainty is a structural limit: no observation map preserving only what remains admissible from the contemporary state can resolve the ambiguity, regardless of how much comparative data is gathered or how sophisticated the reconstruction algorithm becomes. This is not a claim about the practical difficulty of reconstruction but about what information has ceased to exist in the observable projection.

The same consequence applies to any domain in which distinct processes can produce observationally equivalent outcomes. Two musical traditions that arrive at similar harmonic structures, one through independent evolution within a culture's aesthetic constraints and another through cultural contact and adoption, face the same non-identifiability condition if the contemporary harmonic structure determines the same space of admissible continuations regardless of origin. Two gestural systems that develop similar pointing conventions, one through communicative necessity in a specific ecological niche and another through deliberate pedagogical standardization, cannot be distinguished by an observer who has access only to the contemporary gestural repertoire if both produce the same reachable-set structure. In each case, Theorem~4.2 establishes that the observational collision is a property of the projection map rather than a failure of the observer's inferential capacity.

\section{Empirical Hypotheses}

The framework developed in this essay generates several testable predictions concerning cross-modal cognition, though each requires marking explicitly as a hypothesis rather than an established result.

If gesture recognition in human cognition operates as an inverse reconstruction problem attempting to recover latent embodied semantic structure from observed motor trajectories, as discussed in Chapter~\ref{ch:essay-intro}, then the hypothesis is that recognition systems incorporating temporal trajectory information should outperform systems operating solely on isolated frames. The formal motivation is that continuous motor trajectories preserve reachable-set structure that is lost when the trajectory is sampled to discrete configurations. A recognition system that attempts to reconstruct the latent intention by inverting the projection from intention through motor planning to observable configuration should therefore perform better when given access to the continuous trajectory than when given only a sequence of static poses, because the trajectory preserves distinctions that the frame sequence discards. This is testable by comparing recognition accuracy for trajectory-aware versus frame-based algorithms on the same gesture corpus, though such data are not available in the repository's current experiments.

If musical cognition involves reconstructing structured temporal relationships from acoustic projections, then a second hypothesis is that listeners should exhibit systematic reconstruction errors in cases where distinct compositional structures project to similar acoustic surfaces. Specifically, if two different harmonic progressions produce similar spectral and temporal patterns but differ in their reachable continuations—what harmonic moves remain admissible next—then listeners should be more accurate at predicting continuations when given the full compositional structure than when given only the acoustic surface. This predicts measurable differences in continuation judgment tasks depending on whether the task provides structural or surface information, and it further predicts that the error patterns should align with cases where Theorem~4.2 applies: listeners should exhibit greater uncertainty exactly when the observation map produces collisions between distinct compositional histories.

If control processes in biological systems maintain reference values for perceptual variables rather than for output forms, as discussed in Chapter~\ref{ch:essay-pct}, then a third hypothesis is that neural correlates of linguistic behavior should exhibit sustained activity patterns associated with controlled perceptual goals rather than transient patterns associated with individual motor outputs. This predicts that neuroimaging studies should find persistent activation in regions associated with comprehension monitoring, social feedback processing, and effort evaluation throughout utterance production, rather than activation patterns that rise and fall with individual articulatory gestures. The prediction is testable through targeted neuroimaging protocols, though it requires careful experimental design to distinguish sustained goal-related activity from preparatory or anticipatory motor activity, and it remains an empirical hypothesis rather than a consequence of the formal framework.

The distinction between the established consequences in the preceding section and these empirical hypotheses is fundamental. Theorem~4.2 proves that observation maps factoring through reachable-set structure are non-injective under the stated conditions, and this applies to any domain satisfying those conditions, including historical linguistics. The theorem does not prove that human gesture recognition operates as an inverse reconstruction problem, that musical cognition involves reachable-set recovery, or that linguistic behavior is generated by perceptual control loops. Each of these is an empirical conjecture motivated by the formal framework but requiring independent validation through targeted experiments, behavioral measurements, or neurophysiological data. The value of stating them explicitly as hypotheses is that they become testable predictions with clear success and failure conditions rather than remaining at the level of informal speculation, and they illustrate how the admissibility framework can generate empirically consequential predictions across modalities while respecting the boundary between what the mathematics establishes and what the world must be observed to determine.
